% Options for packages loaded elsewhere
\PassOptionsToPackage{unicode}{hyperref}
\PassOptionsToPackage{hyphens}{url}
\documentclass[
]{article}
\usepackage{xcolor}
\usepackage{amsmath,amssymb}
\setcounter{secnumdepth}{-\maxdimen} % remove section numbering
\usepackage{iftex}
\ifPDFTeX
  \usepackage[T1]{fontenc}
  \usepackage[utf8]{inputenc}
  \usepackage{textcomp} % provide euro and other symbols
\else % if luatex or xetex
  \usepackage{unicode-math} % this also loads fontspec
  \defaultfontfeatures{Scale=MatchLowercase}
  \defaultfontfeatures[\rmfamily]{Ligatures=TeX,Scale=1}
\fi
\usepackage{lmodern}
\ifPDFTeX\else
  % xetex/luatex font selection
\fi
% Use upquote if available, for straight quotes in verbatim environments
\IfFileExists{upquote.sty}{\usepackage{upquote}}{}
\IfFileExists{microtype.sty}{% use microtype if available
  \usepackage[]{microtype}
  \UseMicrotypeSet[protrusion]{basicmath} % disable protrusion for tt fonts
}{}
\makeatletter
\@ifundefined{KOMAClassName}{% if non-KOMA class
  \IfFileExists{parskip.sty}{%
    \usepackage{parskip}
  }{% else
    \setlength{\parindent}{0pt}
    \setlength{\parskip}{6pt plus 2pt minus 1pt}}
}{% if KOMA class
  \KOMAoptions{parskip=half}}
\makeatother
\usepackage{longtable,booktabs,array}
\newcounter{none} % for unnumbered tables
\usepackage{calc} % for calculating minipage widths
% Correct order of tables after \paragraph or \subparagraph
\usepackage{etoolbox}
\makeatletter
\patchcmd\longtable{\par}{\if@noskipsec\mbox{}\fi\par}{}{}
\makeatother
% Allow footnotes in longtable head/foot
\IfFileExists{footnotehyper.sty}{\usepackage{footnotehyper}}{\usepackage{footnote}}
\makesavenoteenv{longtable}
\setlength{\emergencystretch}{3em} % prevent overfull lines
\providecommand{\tightlist}{%
  \setlength{\itemsep}{0pt}\setlength{\parskip}{0pt}}
% ===== unicode + compile header for ZIMMERMAN_LAW_OF_GRAVITY =====
% TO COMPILE ON OVERLEAF: upload this .tex + a figures/ folder containing
%   fig1_rar.png  fig2_a0z.png  fig3_btfr.png
% Compile with XeLaTeX (recommended) or pdfLaTeX. Download the PDF and publish
% manually on Zenodo. (Not test-compiled locally; unicode is mapped below.)
\usepackage{amssymb}
\usepackage{newunicodechar}
\providecommand{\textsubscript}[1]{\ensuremath{_{\mbox{\scriptsize #1}}}}
\providecommand{\textonehalf}{\ensuremath{\frac12}}
\newunicodechar{₀}{\ensuremath{_{0}}}
\newunicodechar{₁}{\ensuremath{_{1}}}
\newunicodechar{₃}{\ensuremath{_{3}}}
\newunicodechar{₆}{\ensuremath{_{6}}}
\newunicodechar{ₐ}{\ensuremath{_{a}}}
\newunicodechar{¹}{\ensuremath{^{1}}}
\newunicodechar{²}{\ensuremath{^{2}}}
\newunicodechar{⁰}{\ensuremath{^{0}}}
\newunicodechar{⁴}{\ensuremath{^{4}}}
\newunicodechar{⁵}{\ensuremath{^{5}}}
\newunicodechar{⁻}{\ensuremath{^{-}}}
\newunicodechar{Λ}{\ensuremath{\Lambda}}
\newunicodechar{ρ}{\ensuremath{\rho}}
\newunicodechar{π}{\ensuremath{\pi}}
\newunicodechar{σ}{\ensuremath{\sigma}}
\newunicodechar{Υ}{\ensuremath{\Upsilon}}
\newunicodechar{χ}{\ensuremath{\chi}}
\newunicodechar{Ω}{\ensuremath{\Omega}}
\newunicodechar{Δ}{\ensuremath{\Delta}}
\newunicodechar{γ}{\ensuremath{\gamma}}
\newunicodechar{α}{\ensuremath{\alpha}}
\newunicodechar{µ}{\ensuremath{\mu}}
\newunicodechar{√}{\ensuremath{\surd}}
\newunicodechar{×}{\ensuremath{\times}}
\newunicodechar{−}{\ensuremath{-}}
\newunicodechar{·}{\ensuremath{\cdot}}
\newunicodechar{∝}{\ensuremath{\propto}}
\newunicodechar{≈}{\ensuremath{\approx}}
\newunicodechar{≳}{\ensuremath{\gtrsim}}
\newunicodechar{≥}{\ensuremath{\geq}}
\newunicodechar{⇒}{\ensuremath{\Rightarrow}}
\newunicodechar{→}{\ensuremath{\rightarrow}}
\newunicodechar{↔}{\ensuremath{\leftrightarrow}}
\newunicodechar{⟺}{\ensuremath{\Longleftrightarrow}}
\newunicodechar{⋆}{\ensuremath{\star}}
\newunicodechar{★}{\ensuremath{\star}}
\newunicodechar{✓}{\ensuremath{\checkmark}}
\newunicodechar{✗}{\ensuremath{\times}}
\newunicodechar{◐}{\ensuremath{\circ}}
\newunicodechar{½}{\ensuremath{\frac{1}{2}}}
\newunicodechar{§}{\S{}}
\newunicodechar{Ö}{\"{O}}
\newunicodechar{Ü}{\"{U}}
\newunicodechar{é}{\'{e}}
\newunicodechar{³}{\ensuremath{^{3}}}
\newunicodechar{ä}{\"{a}}
\newunicodechar{Ȳ}{\ensuremath{\bar{Y}}}
\newunicodechar{δ}{\ensuremath{\delta}}
\newunicodechar{θ}{\ensuremath{\theta}}
\newunicodechar{κ}{\ensuremath{\kappa}}
\newunicodechar{λ}{\ensuremath{\lambda}}
\newunicodechar{μ}{\ensuremath{\mu}}
\newunicodechar{φ}{\ensuremath{\varphi}}
\newunicodechar{–}{\textendash{}}
\newunicodechar{—}{\textemdash{}}
\newunicodechar{…}{\ldots{}}
\newunicodechar{′}{'}
\newunicodechar{⁶}{\ensuremath{^{6}}}
\newunicodechar{⁷}{\ensuremath{^{7}}}
\newunicodechar{₂}{\ensuremath{_{2}}}
\newunicodechar{₄}{\ensuremath{_{4}}}
\newunicodechar{ℏ}{\ensuremath{\hbar}}
\newunicodechar{∇}{\ensuremath{\nabla}}
\newunicodechar{≔}{\ensuremath{:=}}
\newunicodechar{≤}{\ensuremath{\leq}}
\newunicodechar{≲}{\ensuremath{\lesssim}}
\newunicodechar{⊙}{\ensuremath{\odot}}
\newunicodechar{Σ}{\ensuremath{\Sigma}}
\newunicodechar{η}{\ensuremath{\eta}}
\newunicodechar{ν}{\ensuremath{\nu}}
\newunicodechar{″}{''}
\newunicodechar{⇔}{\ensuremath{\Leftrightarrow}}
\newunicodechar{±}{\ensuremath{\pm}}
\newunicodechar{¼}{\ensuremath{\tfrac{1}{4}}}
\newunicodechar{β}{\ensuremath{\beta}}
\newunicodechar{∈}{\ensuremath{\in}}
\newunicodechar{∞}{\ensuremath{\infty}}
\newunicodechar{≠}{\ensuremath{\neq}}
% --- v3 erratum build: additional symbol coverage ---
\newunicodechar{ω}{\ensuremath{\omega}}
\newunicodechar{τ}{\ensuremath{\tau}}
\newunicodechar{∫}{\ensuremath{\int}}
\newunicodechar{∂}{\ensuremath{\partial}}
\newunicodechar{≫}{\ensuremath{\gg}}
\newunicodechar{⟨}{\ensuremath{\langle}}
\newunicodechar{⟩}{\ensuremath{\rangle}}
\newunicodechar{⁸}{\ensuremath{^{8}}}
\usepackage{bookmark}
\IfFileExists{xurl.sty}{\usepackage{xurl}}{} % add URL line breaks if available
\urlstyle{same}
\hypersetup{
  hidelinks,
  pdfcreator={LaTeX via pandoc}}

\author{}
\date{}

\begin{document}

\section{The Relational Velocity-Dispersion Spread: A
Modified-Gravity-Impossible Signature of History-Dependent
Inertia}\label{the-relational-velocity-dispersion-spread-a-modified-gravity-impossible-signature-of-history-dependent-inertia}

\textbf{Carl Zimmerman} Briar Creek Tech --- carl@briarcreektech.com

\emph{2026-07-17. This paper supersedes and corrects door D3 of ``No
Pump-Free Corner'' (Zenodo concept DOI
\href{https://doi.org/10.5281/zenodo.21179352}{10.5281/zenodo.21179352});
see the concurrent erratum of that work. Self-cites: MI Field Theory
Results 2026
(\href{https://doi.org/10.5281/zenodo.21403470}{10.5281/zenodo.21403470});
the lensing no-go
(\href{https://doi.org/10.5281/zenodo.21418816}{10.5281/zenodo.21418816}).}

\begin{center}\rule{0.5\linewidth}{0.5pt}\end{center}

\subsection{Abstract}\label{abstract}

In a modified-inertia (MI) reading of the MOND phenomenology, inertia is
a time-nonlocal functional of a body's own worldline: its effective
response depends on its acceleration history through a causal memory
kernel \(K(\Box_u/a_0^2)\). We show that this history-dependence
generates a \emph{relational} signature that no modified-gravity (MG)
theory of the standard class can reproduce in its field sector. For any
theory that (P1) sources a gravitational field \(g(x)\) from the baryons
and (P2) moves tracers on weak-equivalence-principle (WEP) geodesics of
a single metric built from that field, a test body's acceleration at a
fixed cluster-centric position is a function of position alone ---
independent of how it arrived there. Consequently the \emph{equilibrium}
internal velocity dispersion of a self-gravitating subsystem, evaluated
at fixed external field, carries exactly zero spread across infall
history: \(d\sigma_\mathrm{int}/d(\text{history}) \equiv 0\) --- and it
does so \emph{structurally}, because the field carries no
worldline/history label to differentiate. We prove this within the class
\(\{\)QUMOND, AQUAL, AeST/TeVeS, \(f(R)\), local-modified-\(g\}\), for
any \(a_0\), any interpolation function, both coefficient footings,
off-adiabatically, and under retardation. The MG-zero baseline is
Milgrom's MG-virial universality (Milgrom 1983, ApJ \textbf{270}:365;
2014, MNRAS \textbf{437}:2531); the interpolation kernel
\(\nu(y)=\sqrt{1+1/y}\) is Milgrom's (Phys. Lett. A \textbf{253}:273,
1999, Eq. 9). The novelty is (a) the general sourced-field--WEP proof of
the exact-zero off-adiabatically and under retardation, (b) the
demonstration that the discriminant survives a strong-anisotropy control
test and is not re-labeled velocity anisotropy, (c) a designed,
degeneracy-differencing observable \(D(\text{zone})\), and (d) the
framework-specific prediction (the \emph{numbers}, not the
discriminating power).

We carry two coefficient footings on every dimensional number: the
canonical pure-\(\Lambda\) footing
\(a_0=cH_\Lambda/Z=9.36\times10^{-11}\ \mathrm{m\,s^{-2}}\) (with
\(Z=\sqrt{32\pi/3}=5.789\)), and an alternative
\(a_0=1.13\times10^{-10}\ \mathrm{m\,s^{-2}}\)
(\(\rho_\mathrm{total}/cH_0\)). We distinguish two physical channels: a
star-orbit channel within one pressure-supported system (sub-percent,
deep-adiabatic, ELT-gated) and the nearer, larger cluster-member
infall-phase external-field-effect (EFE) channel. For the latter we
design
\(D(\text{zone})=\langle\ln(\sigma_\mathrm{int}/\sigma_\mathrm{bary})\rangle_\mathrm{zone}-\langle\cdot\rangle_\mathrm{ancient}\),
computed within a fixed deprojected external-field bin and tagged by
Rhee et al.~(2017) phase-space infall zones, which differences the
shared radial EFE gradient out as a common mode.

The corrected prediction is two-tiered: (i) in the MG field sector the
fixed-field history spread is \emph{exactly zero}, so any genuinely
field-sector history spread is MG-impossible --- theorem-grade; the
observed-existence claim, however, is confound-contingent, because
ordinary non-equilibrium dynamics (violent relaxation, tidal shocking,
incomplete phase-mixing) produce fixed-field, history-correlated spreads
in Newton+dark-matter and MG alike, and must be separated by design, not
asserted by the theorem. (ii) The leading sign, in the first-infall
pre-pericentre zone, is that first-infall members run \textbf{hotter}
than matched long-resident members --- structural and timescale-free for
any positive-averaging or pure-delay causal kernel, robust across a
\(0/125\) orbit scan; this inverts and supersedes the backwards, dated
pericentre sign-flip of the earlier D3. The magnitude is
\textbf{kernel-hostage}: at the framework-committed memory time it
ranges from \(\sim1.3\%\) (low-pass reading) to \(\sim8\)--\(13\%\)
(pure-delay \(|K|=1\) reading) in the fixed-field
\(\ln(\sigma_\mathrm{int}/\sigma_\mathrm{bary})\) contrast, and is not pinned
by the de Sitter--Unruh foundation. At the low end this sits at or below
the observable's own \(\sim1\)--\(2\%\) projection and
\(\sim2\)--\(8\%\) tidal systematic floors. The signature is
\textbf{MI-class-generic} --- it separates any history-dependent inertia
from modified gravity, but not this framework from Milgrom's linear
no-EFE MI (arXiv:2503.07106), which also spreads --- and it does
\textbf{not} test \(a_0\)'s value or the sign postulate \(s=-1\) (both
postulates; the leading sign rides on \(s=-1\)). This is a
prediction-and-methods paper for a designed, currently-underpowered
discriminator, honest in both directions: it manufactures neither a
detection nor a deficit.

\begin{center}\rule{0.5\linewidth}{0.5pt}\end{center}

\subsection{1. Introduction}\label{introduction}

The MOND regularity --- that galaxy rotation curves and
pressure-supported dispersions follow from the baryons alone through a
single acceleration scale \(a_0\) --- admits two very different
microphysical readings. In \textbf{modified gravity} (MG), the baryons
source a modified gravitational field and matter falls freely in it: the
extra physics lives in the field equation. In \textbf{modified inertia}
(MI), the field equation is unchanged and the extra physics lives in the
response of matter to force --- inertia itself becomes a functional of
the body's motion. These readings are notoriously hard to separate. For
settled, phase-mixed systems they are engineered to coincide on the
radial acceleration relation, and most historically proposed
discriminants (rotation-curve shapes, the baryonic Tully--Fisher
relation, the RAR scatter, weak-lensing profiles) turn out to be either
shared between the two classes or degenerate with the same nuisance
parameters.

The framework studied here is a de Sitter--Unruh modified-inertia
completion in which the acceleration scale is horizon-derived,
\(a_0=cH_\Lambda/Z\) with \(Z=\sqrt{32\pi/3}=5.789\), giving the
canonical value \(a_0=9.36\times10^{-11}\ \mathrm{m\,s^{-2}}\); the
framework's own de Sitter--Unruh interpolation is
\(g_\mathrm{obs}=\sqrt{g_\mathrm{bar}^2+g_\mathrm{bar}\,a_0}\),
i.e.~\(\nu(y)=\sqrt{1+1/y}\) with \(y=g_\mathrm{bar}/a_0\). This
interpolation kernel is \emph{identical} to Milgrom's (Phys. Lett. A
\textbf{253}:273, 1999, Eq. 9; astro-ph/9805346); the framework's
distinctive content is the \(cH_\Lambda/Z\) coefficient and a covariant
MI completion in which the inertial response is a retarded,
time-nonlocal functional of the body's own 4-acceleration through a
kernel \(K(\Box_u/a_0^2)\) with an exact memory time (MI Field Theory
Results 2026, {[}10.5281/zenodo.21403470{]}). Throughout, we carry
\textbf{two coefficient footings} on every dimensional number. (A
``footing'' is the choice of which cosmological rate normalizes \(a_0\):
the pure-\(\Lambda\) dark-energy rate \(cH_\Lambda\), canonical, giving
\(a_0=9.36\times10^{-11}\ \mathrm{m\,s^{-2}}\); versus the total-density
rate \(cH_0\), alternate, giving
\(a_0=1.13\times10^{-10}\ \mathrm{m\,s^{-2}}\).) We state plainly at the
outset, and repeat throughout: \(a_0\)'s value and the response sign
\(s=-1\) are \textbf{postulates} of the framework, not derived here, and
nothing in this paper tests them.

The discriminability question has recently sharpened. A companion
analysis (the lensing no-go, {[}10.5281/zenodo.21418816{]}) closed
gravitational lensing as a discriminator: in the framework's disformal
construction the lensing observables become shared with the MG
(AeST-class) realization, so lensing does not separate MI from MG. That
closure removes one of the few candidate handles and makes the search
for a \emph{genuinely} MG-impossible observable more pressing. This
paper reports one --- with the caveat, developed below, that the
impossibility is airtight in the MG \emph{field sector} and that the
\emph{observable} existence claim must beat same-signed dynamical
confounds by design.

The idea rests on the definitional difference between the two classes.
If inertia is history-dependent, then a self-gravitating subsystem's
internal velocity dispersion depends on the subsystem's acceleration
history --- two otherwise-identical subsystems sitting at the same
place, having arrived by different paths, carry different effective
inertia and therefore different internal heat. If, instead, the extra
physics lives in a sourced field and matter moves on WEP geodesics of
it, then a body's acceleration at a fixed position is fixed by
\emph{where it is}, not \emph{how it got there}, and the field-sector
history spread vanishes identically. The MG-zero side of this statement
is not new in spirit: it is Milgrom's MG-virial universality (Milgrom
1983, ApJ \textbf{270}:365; 2014, MNRAS \textbf{437}:2531), the fact
that in field-based MOND the internal dynamics of a subsystem are set by
its instantaneous configuration and environment. What is new here is (a)
a general proof that the exact-zero holds for the entire
sourced-field--WEP class off-adiabatically and under retardation, not
just for a specific Lagrangian; (b) the demonstration that the resulting
discriminant survives a strong-anisotropy control test and is therefore
not re-labeled velocity anisotropy; (c) a specific designed observable
that differences out the shared, degenerate radial gradient; and (d) the
framework-specific prediction of magnitude and leading sign --- the
numbers, not the discriminating power, which is MI-class-generic.

We are honest in both directions throughout, because both directions are
load-bearing. The MG-impossibility is airtight where it is airtight ---
in the field sector at fixed \emph{true} external field --- and we label
it a theorem there and nowhere else. The detection, by contrast, is
currently out of reach: the magnitude is small and kernel-hostage, the
confounds are same-signed, and the datasets that carry the signal do not
yet exist at the required size. This is a prediction-and-methods paper.
It manufactures neither a detection nor a deficit.

\begin{center}\rule{0.5\linewidth}{0.5pt}\end{center}

\subsection{2. Mechanism and the MG-impossibility
theorem}\label{mechanism-and-the-mg-impossibility-theorem}

\subsubsection{2.1 History-dependent inertia versus instantaneous
EFE}\label{history-dependent-inertia-versus-instantaneous-efe}

In the covariant MI completion, the inertial response of a body is not a
local function of its instantaneous 4-acceleration \(a^\mu\) but a
causal functional of its worldline through the kernel
\(K(\Box_u/a_0^2)\), where \(\Box_u\) is the d'Alembertian along the
worldline. The kernel has an exact, footing-free memory time
(equation-book E10): \[
\tau_\mathrm{mem} = \frac{2c}{a_0} = \frac{2Z}{H_\Lambda}, \qquad \tau_\mathrm{mem} H_\Lambda = 2Z = 11.58 \ \text{(exact, footing-free)},
\] which evaluates to \(\tau_\mathrm{mem}=203\ \mathrm{Gyr}\) (canonical) /
\(168\ \mathrm{Gyr}\) (alt). Physically, the body's effective inertia at
time \(t\) is a memory-weighted functional of the force it has felt over
the preceding \(\sim\tau_\mathrm{mem}\). Because \(K\) is causal, the
memory-felt loading lags the instantaneous loading.

Modified gravity has no such object. In QUMOND, AQUAL, AeST/TeVeS, and
\(f(R)\), the extra physics is a modification of the field equation for
\(g(x)\); a tracer's acceleration is \(a=g(x(t))\), the sourced field
evaluated at the tracer's current position. The field carries no
\(d/dt(\text{worldline})\) label. Two tracers at the same position, with
different histories, feel identical \(g\).

The \textbf{external field effect (EFE)} exists in both classes and is
the cleanest place to see the difference. A subsystem (a galaxy) falling
into a larger system (a cluster) feels the cluster's external field
\(g_\mathrm{ext}\), which loads (suppresses) the subsystem's internal MOND
boost. In MG this loading is \emph{instantaneous}: the internal dynamics
are set by the current \(g_\mathrm{ext}=g_\mathrm{ext}(\text{position})\). In
MI the loading is a memory-weighted functional of the subsystem's
\(g_\mathrm{ext}\)-history along its infall orbit, governed by the
two-frequency subsystem boost \(\theta(y)\),
\(y=\omega_\mathrm{ex}/\omega_\mathrm{in}\) (Milgrom 2022, Phys. Rev.~D
\textbf{106}:064060). At \emph{fixed current} \(g_\mathrm{ext}\), MI
predicts a spread across infall phase; the MG field sector predicts
none.

\subsubsection{\texorpdfstring{2.2 The sourced-field--WEP proof:
\(d\sigma_\mathrm{int}/d(\text{history})=0\)}{2.2 The sourced-field--WEP proof: d\textbackslash sigma\_\{\textbackslash rm int\}/d(\textbackslash text\{history\})=0}}\label{the-sourced-fieldwep-proof-dsigma_rm-intdtexthistory0}

Define the MG class by two premises:

\begin{itemize}
\tightlist
\item
  \textbf{(P1)} the theory \emph{sources a field} \(g(x)\) from the
  baryons (elliptic quasi-static, or hyperbolic/retarded in general);
\item
  \textbf{(P2)} matter tracers are \emph{WEP geodesics} of the single
  Jordan metric built from that field: a tracer's acceleration at event
  \(x\) is a function of \(x\) (and, with retardation, of the
  \emph{source's} past light cone) --- independent of the tracer's own
  orbit, velocity, or acceleration history.
\end{itemize}

Under (P2), the orbit shape, the infall phase
\(y=\omega_\mathrm{ex}/\omega_\mathrm{in}\), and the velocity \(v\) all label
the \emph{tracer}, and they appear nowhere in the equations for the
internal dynamics. Hence

\begin{quote}
\textbf{Theorem (MG field-sector history spread vanishes).} Under
(P1)--(P2), the equilibrium internal velocity dispersion of a
self-gravitating subsystem, evaluated at fixed gravitational field and
fixed internal baryons, is independent of the subsystem's orbital/infall
history: \[
\frac{\partial\sigma_\mathrm{int}}{\partial y}\equiv 0,\qquad \frac{\partial\sigma_\mathrm{int}}{\partial v}\equiv 0,\qquad \frac{\partial\sigma_\mathrm{int}}{\partial(\text{history})}\equiv 0,
\] identically, for any \(a_0\), any interpolation function, both
coefficient footings, elliptic or retarded field equations, and
off-adiabatically. \emph{The zero is structural: under (P2) no
worldline/history variable enters the internal dynamics, so there is no
argument to differentiate. The hypotheses hold for QUMOND, AQUAL,
AeST/TeVeS, \(f(R)\), and any local-modified-\(g\) theory. The sole
evasion --- a coupling to the tracer's own velocity/worldline ---
violates (P2) and is itself modified inertia.}
\end{quote}

This is verified symbolically with an arbitrary interpolation \(\mu\)
and arbitrary \(a_0\), and numerically across \(\{\)canonical,
alt\(\}\times\{\)framework \(\nu\), standard-MOND \(\nu\),
exponential-RAR\(\}\): four orbit families matched at the same position
return an identical internal boost, with peak-to-peak spread \(=0\)
exactly. We stress that the symbolic zero is \emph{structural}, not a
surprising computation --- the algebraic-MOND EFE simply contains no
worldline/history variable to differentiate --- and that structural
absence is precisely the physical content: a subsystem's internal heat
is fixed by \emph{where it is}, not \emph{how it got there}. This is
Milgrom's MG-virial universality, here stated as an off-adiabatic
theorem for the whole class.

Two structural checks confirm the exact-zero survives the natural
loopholes:

\begin{itemize}
\tightlist
\item
  \textbf{Time-varying potential (off-adiabatic).} Going non-adiabatic
  replaces the fixed position by the worldline \(x(t)\) but attaches no
  velocity or history label to \(g\). A settled member and a deep
  plunger that reach the same current radius by different histories have
  identical MG boost. The MG field is memoryless; the infall history
  drops out exactly.
\item
  \textbf{Retardation / finite crossing.} The field lag is
  \(\sim v/c\sim3.5\times10^{-3}\); retardation lives in the
  \emph{source's} past light cone and is felt identically by every
  member at \((x,t)\). It shifts the \emph{mean} field at the
  \(10^{-3}\) level but adds \emph{zero} infall-phase family spread.
  (Contrast MI, whose kernel is retarded along the \emph{tracer's own}
  worldline.)
\end{itemize}

MI violates (P2) by construction: the inertial response is a functional
of the body's own worldline through \(K(\Box_u/a_0^2)\). Two members at
the same position with different orbital histories carry different
effective inertia, hence a spread. That is the entire field-sector
distinction.

\textbf{Scope of the theorem, stated precisely and up front.} The
exact-zero is a statement about the MG \emph{field channel} at fixed
\emph{true} 3D external field. It is \emph{not} a claim that the total
observed dispersion spread of a real, recently-infalling subsystem is
zero in an MG (or a Newton+dark-matter) universe. Ordinary
non-equilibrium dynamics --- violent relaxation, tidal shocking,
incomplete phase-mixing, ram-pressure, substructure --- produce
fixed-field, history-correlated dispersion spreads that are entirely
standard and are \emph{not} modified inertia. They are shared dynamical
confounds, same-signed with the MI excess, and they are separated (Sec.
4) by radial profile and sign, not by the theorem. The theorem-grade
badge attaches to ``the MG field-sector contribution to the fixed-field
history spread is exactly zero,'' not to ``any measured fixed-field
spread is modified inertia.'' The observable existence claim of Sec. 5.1
is therefore confound-contingent --- it depends on the F3 radial-profile
separator holding --- and is \emph{not} itself theorem-grade.

\subsubsection{2.3 The anisotropy control test (star-orbit
channel)}\label{the-anisotropy-control-test-star-orbit-channel}

The obvious worry is that a ``spread across orbit families'' is just
velocity anisotropy \(\beta(r)\) in disguise --- the classic
mass-anisotropy degeneracy, which lives in the angular sector of the
velocity ellipsoid and is shared by Newton+dark matter, MG, and MI
alike. If the MI signal were re-labeled anisotropy it would carry no
discriminating power.

We tested this in the \emph{star-orbit} channel (Sec. 3.1), where the
relevant statistic is the orbit-family enclosed-mass split \(d\ln M\).
We built a \emph{proper steady-state} distribution function with
strongly radially-varying anisotropy --- \(\beta\) running from
\(+0.13\) in the center to \(-0.72\) outward, i.e.~\(\Delta\beta\sim9\)
between subsamples, with anisotropy a genuine function of the orbital
integrals --- and \emph{no MI}. The mechanism is that MI multiplies the
radial and tangential velocity components \((v_r,v_t)\) by the
\emph{same} per-orbit factor \(f(e)\), which \textbf{cancels in the
anisotropy ratio}
\(\beta=1-\langle v_t^2\rangle/(2\langle v_r^2\rangle)\) and surfaces
only in the enclosed-mass normalization \(M\propto f^2\sigma_r^2\). MI
is orthogonal to \(\beta\) by construction: anisotropy is a statement
about velocity \emph{direction}; MI acts on velocity \emph{magnitude}.

We report the control test honestly, including its demotions from the
verify lane. The pure-anisotropy DF returns a zero-point
\(d\ln M\approx+0.033\), which is \textbf{larger in magnitude than the
MI signal itself} (\(d\ln M_\mathrm{MI}\approx-0.02\)) and of the
\emph{opposite sign}; across DF shapes the zero-point drifts
\(+0.009\to+0.062\). In other words, valid equilibrium anisotropy is
separated from MI by \emph{sign}, not by amplitude: the anisotropy
contamination is signal-sized, and a real measurement would require a
forward Schwarzschild / made-to-measure calibration whose residual is
demonstrated sub-signal --- which we have not done. What the test
\emph{does} establish is that strong, radially-varying, physically-valid
anisotropy does not reach the MI-signed value and cannot fake the MI
sign. It is a sign-immunity result, not a clean amplitude pass.

A separate and sharper practical mimic surfaced in the same lane: a
radius-correlated \emph{non-steady-state population mix} (a subsystem
out of equilibrium, with the mix varying with radius) drives
\(d\ln M\approx-0.71\), roughly \(35\times\) the MI signal and
\emph{MI-signed}. This is the dominant unguarded false-detection channel
for the star-orbit statistic, and it is not beaten by the anisotropy
argument --- it is beaten only by an equilibrium/relaxed-population
selection. We carry it forward into the confound budget (Sec. 4.4,
F-table) rather than treating the discriminant as clean.

\subsubsection{2.4 The one evasion is definitionally modified
inertia}\label{the-one-evasion-is-definitionally-modified-inertia}

Is there \emph{any} force that manufactures the orbit-family spread
while remaining modified gravity? We stress-tested the boundary of the
class:

\begin{itemize}
\tightlist
\item
  \textbf{Gravitomagnetic} \(F=m\,v\times B_g\): velocity-dependent but
  antisymmetric, \(F\cdot v\equiv0\), does no work, cannot heat a member
  --- zero spread.
\item
  \textbf{Dissipative drag} \(F=-\gamma(x)v\): non-conservative, no
  medium in a collisionless system; if imposed it \emph{cools} toward
  \(\sigma\to0\), no steady spread; not a field theory of \(g\).
\item
  \textbf{\(f(R)\) / chameleon} \(G_\mathrm{eff}(x)\):
  environment-dependent coupling is a function of local
  density/potential, hence of \(x\), shared by all tracers --- spread
  zero.
\item
  \textbf{Retarded MG}: retardation is in the source, felt identically
  by all members at \((x,t)\) --- spread zero.
\item
  \textbf{Disformal / Finsler / SME velocity-dependent coupling}
  \(\delta a\sim\beta_c(v\cdot\nabla\phi)\): this \emph{does} create an
  orbit split (e.g.~\(\sim9.5\%\) at \(\beta_c=0.1\)). But it couples to
  the \emph{tracer's own velocity}, breaking (P2) and WEP --- it is
  modified \textbf{inertia} in an MG costume. Standard MG is WEP-exact
  (\(\beta_c\equiv0\)).
\end{itemize}

The only door that opens is not an MG door. Any theory that manufactures
a field-sector orbit-family spread has, by construction, made the
tracer's acceleration depend on how it moves --- it has put history into
the \emph{inertial} response. Whatever one names it, a finite
field-sector fixed-field history spread \textbf{is} the modified-inertia
signal. There is no pure sources-a-field-\(g(x)\) channel to it.

\begin{center}\rule{0.5\linewidth}{0.5pt}\end{center}

\subsection{3. The two channels}\label{the-two-channels}

The history spread appears in two physically distinct places, which the
banked analysis initially conflated and a session correction separated.
They differ by roughly an order of magnitude in amplitude and are
measured by entirely different means.

\subsubsection{3.1 Channel A --- star-orbit spread within one
pressure-supported
system}\label{channel-a-star-orbit-spread-within-one-pressure-supported-system}

Inside a single dSph or elliptical, every star at radius \(r\) feels the
same \(g_\mathrm{bar}(r)\). A \emph{local} \(\mu(|a|)\) inertia would give
every star the same \(|a|\) and no spread. The MI effect is
non-adiabatic: a star on an \emph{eccentric} orbit time-samples a
\emph{varying} \(|a|\) (large at pericenter), so its memory-averaged
effective inertia differs from that of a circular orbit at the same
energy --- a Jensen gap over the curvature of the nonlinear \(\nu\).
Different orbital families at the same radius carry different effective
inertia, hence an intrinsic dispersion spread.

The magnitude is set by \(\tau_\mathrm{mem}\) versus the orbital time
\(T_\mathrm{orb}\approx2\pi r/\sigma\). For every real pressure-supported
system this ratio is enormous:

{\def\LTcaptype{none} % do not increment counter
\begin{longtable}[]{@{}
  >{\raggedright\arraybackslash}p{(\linewidth - 4\tabcolsep) * \real{0.3333}}
  >{\raggedright\arraybackslash}p{(\linewidth - 4\tabcolsep) * \real{0.3333}}
  >{\raggedright\arraybackslash}p{(\linewidth - 4\tabcolsep) * \real{0.3333}}@{}}
\toprule\noalign{}
\begin{minipage}[b]{\linewidth}\raggedright
system
\end{minipage} & \begin{minipage}[b]{\linewidth}\raggedright
\(\tau_\mathrm{mem}/T_\mathrm{orb}\) (canonical)
\end{minipage} & \begin{minipage}[b]{\linewidth}\raggedright
regime
\end{minipage} \\
\midrule\noalign{}
\endhead
\bottomrule\noalign{}
\endlastfoot
Draco / Sculptor / Fornax dSph & \(1367\times\) / \(1086\times\) /
\(544\times\) & deep adiabatic \\
Crater II (diffuse) & \(83\times\) & deep adiabatic \\
Coma cluster members & \(22\times\) & deep adiabatic \\
\end{longtable}
}

Because \(\tau_\mathrm{mem}=203/168\ \mathrm{Gyr}\gg T_\mathrm{orb}\), every
orbit sits far above the kernel's edge frequency
\(a_0/2c=1/\tau_\mathrm{mem}\). On the pure-phase branch the memory
magnitude \emph{saturates and freezes} at the orbit-mean pre-history
fixed point --- there is \textbf{no resonant amplification}. The spread
is the small residual adiabatic Jensen gap. Direct orbit integration
through the real memory kernel (an independent, 19/19-verified
integrator) gives an eccentric-orbit RAR offset \(<0.007\) dex out to
\(e\approx0.9\), sign negative, and an isotropic-dSph ensemble
\(\nu_\mathrm{eff}/\nu_\mathrm{circ}=0.990\)--\(0.997\). Over realistic
eccentricity distributions the RMS relational \(\sigma\)-spread is \[
\text{Channel A: } \sim0.2\text{–}0.35\%\ \text{(fiducial cored dSph)},\quad \lesssim1\%\ \text{(point-mass ceiling)},\quad <0.1\%\ \text{(ellipticals, } y\gg1).
\] Both footings agree to \(<20\%\) (\(a_0\) cancels at fixed depth
\(y\)). The sign is negative: eccentric orbits present a slightly lower
effective \(\nu\) and run marginally cooler. This channel is real but
small --- an order of magnitude below the cluster channel --- and, as
Sec. 6 shows, needs ELT-class per-star velocities plus a per-star 3D
orbit tag that does not yet exist.

\subsubsection{3.2 Channel B --- cluster-member infall-phase EFE
spread}\label{channel-b-cluster-member-infall-phase-efe-spread}

The nearer, larger channel is the \emph{subsystem} EFE. A whole
cluster-member galaxy's internal dispersion is loaded by the cluster
external field through the two-frequency boost \(\theta(y)\),
\(y=\omega_\mathrm{ex}/\omega_\mathrm{in}\) (Milgrom 2022). Here \(y\) is the
member's dimensionless internal depth: small \(y\) means the member sits
deep in the modified regime (a diffuse, low-internal-acceleration
galaxy), and \emph{only} diffuse low-\(y\) members carry the signal. It
is essential to separate two quantities that the banked notes stress
repeatedly:

{\def\LTcaptype{none} % do not increment counter
\begin{longtable}[]{@{}
  >{\raggedright\arraybackslash}p{(\linewidth - 4\tabcolsep) * \real{0.3333}}
  >{\raggedright\arraybackslash}p{(\linewidth - 4\tabcolsep) * \real{0.3333}}
  >{\raggedright\arraybackslash}p{(\linewidth - 4\tabcolsep) * \real{0.3333}}@{}}
\toprule\noalign{}
\begin{minipage}[b]{\linewidth}\raggedright
\end{minipage} & \begin{minipage}[b]{\linewidth}\raggedright
current-configuration \(\theta(y_\mathrm{cur})\) boost
\end{minipage} & \begin{minipage}[b]{\linewidth}\raggedright
MG-impossible history spread
\end{minipage} \\
\midrule\noalign{}
\endhead
\bottomrule\noalign{}
\endlastfoot
\textbf{magnitude} & \(6\)--\(13\%\) & \(\sim1.3\%\) (low-pass) to
\(\sim8\)--\(13\%\) (pure-delay) \\
\textbf{status} & partly MG-\textbf{shared} (MG has its own
instantaneous EFE) & MG-\textbf{impossible} in the field sector \\
\textbf{is it the discriminant?} & \textbf{NO} --- absorbed by an
\(a_0\)-rescale / EFE term & \textbf{YES} --- this is the signal \\
\end{longtable}
}

\begin{itemize}
\tightlist
\item
  \textbf{The instantaneous \(\theta(y_\mathrm{cur})\) boost} --- the
  \emph{current-configuration} EFE contrast across infall phase. For a
  diffuse deep-MOND member (\(a_\mathrm{in}\approx0.3\,a_0\) internal,
  \(a_\mathrm{ex}\approx a_0\) external, at the transition shell) this is
  the banked \textbf{6--13\%} band (kernel corners \(\theta(0)=\sqrt2\)
  floor \(\to5.5\%\), \(\theta(0)=2\) fiducial \(\to9.5\%\),
  \(\theta(0)=e\) ceiling \(\to11.5\%\) over \(y\le1.5\); fiducial
  reproduction \(9.5\%\), both footings). This is a
  current-configuration quantity that MG partly shares --- MG has its
  own instantaneous \(\theta_\mathrm{MG}(y_\mathrm{cur})\) EFE --- and it is
  therefore \emph{not} the discriminant. A constant EFE boost is
  absorbed by an \(a_0\)-rescale or an EFE term.
\item
  \textbf{The MG-impossible history spread} --- the residual spread
  \emph{at fixed current field and fixed \(y_\mathrm{cur}\)}, riding on top
  of that boost, sourced only by the memory-weighted difference between
  the felt and the current external field. This is the discriminant. Its
  magnitude is set by \(\tau_\mathrm{mem}\) \emph{and by the kernel shape},
  and here the framework's own foundation is genuinely underdetermined
  (Sec. 5.4): under a decaying \textbf{low-pass} reading of the
  committed E10 memory, both members lag the field nearly equally and
  the residual is \emph{residence-time-limited} to \(\sim1.3\%\)
  (canonical) / \(\sim1.5\%\) (alt); under the framework's committed
  \textbf{pure-phase / all-pass} kernel (\(|K|=1\), equation-book E13),
  the felt field is the current field delayed by an uncomputed group
  delay, and a resolvable transient survives, taking the fixed-field
  first-infall-vs-ancient contrast up to \(\sim8\)--\(13\%\). Both
  readings share the \emph{sign} (Sec. 5.2); only the magnitude differs.
  We write the committed-memory magnitude as a band: \[
  \text{Channel B (MG-impossible piece): } \sim1.3\text{–}13\%\ \text{in the fixed-field } \ln(\sigma_\mathrm{int}/\sigma_\mathrm{bary})\ \text{contrast (both footings, kernel-shape band)}.
  \]
\end{itemize}

The magnitude is \textbf{kernel-hostage}, i.e.~it depends on the
undetermined loading function \(\theta(y)\) and the kernel shape and is
quoted as a band, not a derived number: \(\theta(y)\) is not derived by
the de Sitter--Unruh foundation, which fixes only the cone endpoints.
The existence of the spread (in the field sector), its leading sign
(Sec. 5), and MG \(=0\) are the load-bearing claims; the amplitude is a
band.

Two structural facts about Channel B matter for the design. First, the
fractional amplitude at fixed dimensionless depth is
\emph{mass-independent} --- cluster mass sets the transition-shell
radius \(R(a_\mathrm{ex}=a_0)=\sqrt{GM/a_0}\) (\(0.39\ \mathrm{Mpc}\) at
\(10^{14}\), \(1.22\ \mathrm{Mpc}\) at \(10^{15}\,M_\odot\), canonical)
and the crossing/memory window, but not the amplitude. Second, only
\emph{diffuse} members (low \(\omega_\mathrm{in}\), low \(y\)) reach the
signal-bearing regime: UDGs (\(\sigma\sim15\)) and dSph-scale members
(\(\sigma\sim10\)--\(50\)) are the carriers, whereas \(L^\ast\)
ellipticals (\(\sigma\sim200\)) are internally near-Newtonian
(\(y\gg1\)) and adiabatically dead. The survey-bright members are
exactly the dead ones --- which, as Sec. 6 shows, is the power wall.

\begin{center}\rule{0.5\linewidth}{0.5pt}\end{center}

\subsection{\texorpdfstring{4. The observable
\(D(\text{zone})\)}{4. The observable D(\textbackslash text\{zone\})}}\label{the-observable-dtextzone}

\subsubsection{4.1 The key degeneracy and the design
principle}\label{the-key-degeneracy-and-the-design-principle}

Both MI and MG have an EFE that loads the internal boost with the
\emph{current} \(g_\mathrm{ext}\). So in \emph{both} theories
\(\sigma_\mathrm{int}/\sigma_\mathrm{bary}\) varies with cluster-centric
radius --- the shared \textbf{radial EFE gradient}. That radial trend is
not the MI signal; it is the killer common mode. The distinctive MI
signal is the \emph{residual spread at fixed cluster-centric radius
(fixed current field), correlated with infall history} --- for which the
MG field sector predicts exactly zero.

The design principle is to project
\(\sigma_\mathrm{int}/\sigma_\mathrm{bary}\) onto the (radius \(\times\)
phase) plane and difference \emph{along the phase axis within a
fixed-radius bin}. The shared radial gradient is a common mode on the
radius axis with no phase label; it cancels to the bin-width residual
(the projection alias of Sec. 4.4). The MG field-sector residual at
fixed \emph{true} radius is exactly zero (the theorem); only the MI
history signal --- and a controllable projection alias --- survive.

\subsubsection{4.2 The statistic}\label{the-statistic}

Within one deprojected external-field bin (with \(a_\mathrm{ext}\) from the
caustic mass profile, width \(\le0.3\) dex): \[
D(\text{zone})\equiv\big\langle\ln[\sigma_\mathrm{int}/\sigma_\mathrm{bary}]\big\rangle_\mathrm{zone}-\big\langle\ln[\sigma_\mathrm{int}/\sigma_\mathrm{bary}]\big\rangle_\mathrm{ancient},
\] with the infall-phase proxy = the Rhee et al.~(2017, ApJ
\textbf{843}:128) projected-phase-space zone read off from
\((R_\mathrm{proj}/r_{200},\ |v_\mathrm{los}-v_\mathrm{cl}|/\sigma_\mathrm{cl})\),
and \(\sigma_\mathrm{bary}\) from the anisotropy-immune Wolf et al.~(2010,
MNRAS \textbf{406}:1220) half-light mass. The four Rhee zones, with
their physical meaning, are: \textbf{ancient-infall} (long-virialized,
the reference), \textbf{first-infall} (falling in for the first time,
pre-first-pericentre), \textbf{recent-infall} (just past first
pericentre), and \textbf{backsplash} (out again, past first apocentre).
Operationally, \(D(\text{zone})\) is the quantity an IFU delivers
directly: the difference in
\(\langle\ln(\sigma_\mathrm{int}/\sigma_\mathrm{bary})\rangle\) between a
phase zone and the ancient reference, at matched external field.

\textbf{MG field sector: \(D(\text{zone})=0\) for every zone at fixed
true radius} (symbolic \(d/dy=0\), any \(a_0\), any interpolation, both
footings). \textbf{MI: a nonzero, history-correlated pattern}, whose
leading term is the first-infall excess of Sec. 5. The measured
\(D(\text{first-infall})\) at the framework-committed memory is
\(\approx+1.3\%\) (canonical) / \(\approx+1.5\%\) (alt) under the
low-pass kernel reading --- the difference between the first-infall zone
(\(+0.81\%\) dev-vs-mean, canonical) and the ancient zone (\(-0.47\%\))
--- and up to \(\sim8\)--\(13\%\) under the pure-delay reading (Sec.
5.4). This is the number an observer's
\(\langle\ln\sigma_\mathrm{int}/\sigma_\mathrm{bary}\rangle\) contrast equals,
and it is the number that must be compared to the systematic floors
below.

\subsubsection{4.3 Anisotropy immunity via the Wolf
normalization}\label{anisotropy-immunity-via-the-wolf-normalization}

The member's own internal anisotropy \(\beta\) enters
\(\sigma_\mathrm{los}\), and infall can induce radial \(\beta\),
threatening a phase alias. The Wolf mass
\(M(r_{1/2})=3\langle\sigma_\mathrm{los}^2\rangle r_{1/2}/G\) is
\(\beta\)-immune to first order at the half-light radius. The residual
Wolf \(\beta\)-leak is \(\lesssim1.8\%\) and \emph{monotone} in
\(\beta\) --- it cannot produce a sign-flip and folds into the
same-signed heating-only confound family. This is a \emph{softer}
control than the star-orbit sign-immunity of Sec. 2.3 (it is a
first-order cancellation, not an orthogonality); with IFU 3D internal
kinematics \(\beta\) is measured directly and the leak is calibrated.

\subsubsection{4.4 Beating the projection alias and the same-signed
confounds}\label{beating-the-projection-alias-and-the-same-signed-confounds}

The one MG evasion of \(D(\text{zone})\) is \textbf{projection}: at
fixed \emph{projected} radius, radial plungers sit at a different
\emph{true} radius than settled members, so MG's real radial trend
aliases into the phase axis. This is honestly the sharpest limit of the
observable. An isotropic Monte Carlo puts the raw alias at
\(\sim2.2\)--\(2.4\%\); a class-\emph{blind} scalar deprojection barely
helps (\(2.25\%\to2.19\%\)), because the alias is driven by a
class-\emph{dependent} LOS-depth residual, not a mean bias. Two
consequences follow, both toward caution:

\begin{itemize}
\tightlist
\item
  The Rhee zones \emph{are} the orbit-class-aware deprojection --- they
  are calibrated on N-body orbital history, not on a scalar radial mean
  --- which is exactly the class-aware correction a scalar deprojection
  cannot supply. But the tagging has finite purity \(p<1\); at
  \(p=0.5\)--\(0.9\) the residual alias is \(\sim1.3\)--\(2.0\%\), not
  the \(\sim0.01\%\) one would get by binning on the (unobservable)
  simulation-true radius. The mitigation chain (caustic membership +
  Dressler--Shectman substructure cut + class-aware zone deprojection +
  relaxed-cluster selection) is \textbf{load-bearing, not optional}: a
  ``detection'' that bins by projected radius and skips the cuts
  measures projection and interlopers, not modified inertia.
\item
  Anisotropic / filamentary infall (plungers entering along the major
  axis / LOS) can push the raw alias non-monotonically to \(\sim7\%\)
  --- band-sized --- and biases the spherical caustic
  \(a_\mathrm{ext}(R)\). Quantifying this needs a triaxial-potential,
  infall-axis forward model; the banked isotropic figure is a lower
  bound.
\end{itemize}

\textbf{The signal-to-systematic verdict, at the committed memory.} This
is the honest crux. Under the low-pass kernel reading the measurable
\(D(\text{first-infall})\approx1.3\)--\(1.5\%\) absolute sits \emph{at
or below} its own systematic floors: the projection alias is
\(\sim1.3\)--\(2.0\%\) (up to \(\sim7\%\) filamentary/triaxial) and the
same-signed tidal floor is \(\sim2\)--\(8\%\). In that corner the signal
is comparable to, or below, the floors, and the honest expectation is a
non-detection (Sec. 6.2). The signal clears its floors only under the
pure-delay kernel corner (up to \(\sim8\)--\(13\%\)) --- which is a
genuinely uncomputed reading of the same committed memory, not a safe
baseline. The exploratory S/N figures quoted in the scout (Sec. 6.2)
that reach \(\sim2\)--\(2.5\sigma\) used a signal of \(\sim9.5\%\) (the
top of the band); at the low-pass \(\sim1.3\%\) the honest verdict
slides to the scout's central expectation, a non-detection at every
accessible sample size.

The same-signed dynamical confounds (tidal heating/stripping,
ram-pressure, environmental quenching, non-equilibrium population mix)
are beaten not by the phase-difference alone but by a joint four-part
fingerprint that only MI trips:

{\def\LTcaptype{none} % do not increment counter
\begin{longtable}[]{@{}
  >{\raggedright\arraybackslash}p{(\linewidth - 8\tabcolsep) * \real{0.2000}}
  >{\raggedright\arraybackslash}p{(\linewidth - 8\tabcolsep) * \real{0.2000}}
  >{\raggedright\arraybackslash}p{(\linewidth - 8\tabcolsep) * \real{0.2000}}
  >{\raggedright\arraybackslash}p{(\linewidth - 8\tabcolsep) * \real{0.2000}}
  >{\raggedright\arraybackslash}p{(\linewidth - 8\tabcolsep) * \real{0.2000}}@{}}
\toprule\noalign{}
\begin{minipage}[b]{\linewidth}\raggedright
source
\end{minipage} & \begin{minipage}[b]{\linewidth}\raggedright
F1 fixed-\(r\) phase contrast
\end{minipage} & \begin{minipage}[b]{\linewidth}\raggedright
F2 leading sign (first-infall hotter)
\end{minipage} & \begin{minipage}[b]{\linewidth}\raggedright
F3 rises outward
\end{minipage} & \begin{minipage}[b]{\linewidth}\raggedright
F4 baryon-blind
\end{minipage} \\
\midrule\noalign{}
\endhead
\bottomrule\noalign{}
\endlastfoot
\textbf{MI (this framework, MI-class)} & ✓ & ✓ & ✓ & ✓ \\
MG (true \(r\)) & ✗ (\(=0\)) & ✗ & ✗ & ✓ \\
MG projection alias & ✓ & ✗ & ✗ & killed by zone deprojection \\
interlopers (uncut) & ✓ & ✗ & ✗ & killed by DS + caustic \\
tidal heating/stripping & ✓ & ✗ & ✗ (inward) & ✗ (tidal features) \\
ram-pressure / quenching & ✓ & ✗ & ✗ & ✗ (gas/SF marks) \\
non-equilibrium population mix & ✓ & can be MI-signed & ✗ (not
outward-peaked) & ✗ (mixed populations) \\
\end{longtable}
}

The \textbf{radial-profile separator (F3)} is the most robust: MI rises
\emph{outward}, peaking at the MOND-transition shell
\(a_\mathrm{ext}\sim0.3\)--\(1\,a_0\) (\(\sim R_{500}\)--\(R_{200}\)) and
dying in the core, whereas tidal heating rises \emph{inward} toward
pericenter --- opposite slopes. The \textbf{baryon-blind split (F4)} ---
environmental confounds mark the baryons (gas stripping, burst/truncated
SF, tidal morphology) while the inertia signal does not --- is
corroborating, not decisive: it degrades for a \emph{dry} tidal-heating
episode on a gas-poor dE, exactly the diffuse carriers that carry the
signal. The non-equilibrium population mix (Sec. 2.3,
\(d\ln M\approx-0.71\), MI-signed) is the most dangerous mimic because
it can carry the MI sign; it is separated by F3 (it is not
outward-peaked at the transition shell) and by relaxed-cluster /
equilibrium-member selection. Treat F3 as the primary separator, sign
(F2) as necessary-not-sufficient, and F4 as support. The
observable-existence claim of Sec. 5.1 is theorem-grade \emph{only
through} this separation holding; the theorem itself is the field-sector
zero.

\begin{center}\rule{0.5\linewidth}{0.5pt}\end{center}

\subsection{5. The prediction}\label{the-prediction}

The prediction is two-tiered, and it \textbf{supersedes and corrects}
the earlier D3 pre-registration (the dated pericentre sign-flip of ``No
Pump-Free Corner,'' {[}10.5281/zenodo.21179352{]}; see the concurrent
erratum).

\subsubsection{5.1 Tier (i): existence --- theorem-grade in the field
sector, confound-contingent as an
observable}\label{tier-i-existence-theorem-grade-in-the-field-sector-confound-contingent-as-an-observable}

\begin{quote}
At fixed cluster-centric gravitational field, a modified-inertia
universe exhibits a nonzero internal-\(\sigma\) spread correlated with
infall history. In the field sector of any instantaneous-EFE gravity
(QUMOND/AeST/\(f(R)\); Milgrom's MG-virial universality) this
fixed-field history spread is exactly zero.
\end{quote}

The theorem-grade content is the field-sector zero: \emph{no
sourced-field--WEP theory contributes any fixed-field history spread}.
It is independent of the sign and of the \(s=-1\) postulate, and it
holds for any \(a_0\), any interpolation, both footings,
off-adiabatically, and under retardation. The \emph{observable}
existence claim --- that a measured nonzero fixed-field spread evidences
modified inertia --- is one step weaker: it is confound-contingent,
because non-equilibrium dynamics, tidal shocking, and projection produce
same-signed fixed-field spreads in Newton+dark-matter and MG alike (Sec.
2.2, 2.3, 4.4). What makes a \emph{measured} spread MG-impossible is not
the raw existence but the F3 outward-rising radial profile peaking at
the transition shell, which the confounds do not share. We state this
distinction as the honest core of the paper: the impossibility is
airtight for the field sector and rides on the confound separation for
the observable.

\subsubsection{5.2 Tier (ii): leading sign --- first-infall members run
hotter}\label{tier-ii-leading-sign-first-infall-members-run-hotter}

\begin{quote}
Among members matched at the same cluster-centric field,
\textbf{first-infall pre-pericentre members are hotter} (larger internal
\(\sigma\)) than matched long-resident / post-pericentre members;
equivalently, the \(\sigma\)-excess decreases monotonically with
accumulated loading (\(\approx\) time-since-infall).
\end{quote}

This sign is \textbf{structural and timescale-free} for the physical
class of kernels. On a monotonically \emph{rising} field approach from a
low-field past, the causal memory-felt external field is \emph{always
below} the current field, so the member is under-loaded, hence
\(\nu\)-boosted hotter --- for \emph{any positive-averaging or
pure-delay} causal kernel. (The one caveat: a fully dispersive all-pass
kernel with a signed impulse response could in principle
\emph{overshoot} a rising ramp near onset, felt \(>\) current, flipping
the sign locally; the framework's committed \(|K|=1\) branch was checked
in the pure-delay limit and preserves the sign, and the physical reading
is scoped to the monotonic pre-pericentre approach where overshoot does
not occur.) The result is robust across the full grid --- memory time
\(0.1\)--\(203\ \mathrm{Gyr}\), both kernel shapes (exponential low-pass
and the pure-phase group-delay branch), member depth
\(0.1\)--\(1.0\,a_0\), pre-infall field \(0\)--\(0.3\,a_0\), both
footings --- and across an independent 125-point orbit-distribution scan
(masses \(10^{14}\)--\(5\times10^{14}\,M_\odot\), apocenter
\(2\)--\(5\ \mathrm{Mpc}\), pericenter \(0.2\)--\(0.6\ \mathrm{Mpc}\)):
\textbf{0/125 sign flips}. The signal is essentially monotone in
accumulated loading, not a sharp dated event.

The corrected falsifier is therefore inverted relative to the published
D3: a significantly \textbf{negative} fixed-field sign --- first-infall
members measured \emph{cooler} at \(\ge3\sigma\) --- or a null spread at
adequate power, is what would falsify Tier (ii). (What the published D3
named as the \emph{signature} --- a coherent pre-pericentre deficit ---
is in fact the \emph{falsifier}.) The sign is conditional on the
\(s=-1\) postulate; \(s=+1\) reverses it.

\subsubsection{5.3 What is retracted, and
why}\label{what-is-retracted-and-why}

The earlier D3 pre-registered a \textbf{dated pericentre sign-flip}:
first-infall pre-pericentre members \(-11\) to \(-21\%\) \emph{cold},
flipping to a post-pericentre \emph{excess}, decisive at
\(3\sigma\approx2029\)--\(2031\). Both the polarity and the timescale
were wrong.

\begin{enumerate}
\def\labelenumi{\arabic{enumi}.}
\item
  \textbf{Polarity inversion (a bookkeeping error).} The ``cold isolated
  past'' of a first-infall member was encoded as a \emph{low} value of
  the loading ratio \(y=\omega_\mathrm{ex}/\omega_\mathrm{in}\). But the
  loading factor \(\theta(y)\) is \emph{maximal} as \(y\to0\); true
  isolation is \(a_\mathrm{ext}\to0\) (\emph{zero} external loading for any
  \(\theta\)), not low-\(y\). Holding a fixed nonzero \(a_\mathrm{ext}\)
  while sending \(y_\mathrm{hist}\to0\) injects \emph{maximal} past loading
  and produces a spurious deficit. Worked correctly in field space, a
  first-infall member on a rising-field approach has a memory-felt field
  \emph{below} its current field --- under-loaded, hence \emph{hotter,
  not colder}. A parallel text-label slip in a companion script (echoed
  in the banked GAP\_STATEMENT E4/E7) printed ``plungers less boosted''
  while its own loop output them hotter --- the same conflation of ``low
  \(\theta\)'' with ``low boost,'' when low \(\theta\) means \emph{less}
  suppression and therefore \emph{more} boost. The claimed
  raw-loading-versus-memory ``competition'' was an artifact of this
  encoding; in field space the two contributions reinforce, and the net
  sign is fixed unambiguously by
  \(\mathrm{sign}(a_\mathrm{ext,felt}-a_\mathrm{ext,now})\).
\item
  \textbf{Timescale hostage (a physics error).} The dated flip was
  computed with a Lorentzian memory of \(\tau=0.45\ \mathrm{Gyr}\) (a
  dwarf-sector value), which is not the framework's committed covariant
  memory. The equation-book kernel gives
  \(\tau_\mathrm{mem}=2c/a_0=2Z/H_\Lambda=203/168\ \mathrm{Gyr}\),
  footing-free and algebraic from the kernel. Against a cluster crossing
  time of \(\sim1\)--\(2\ \mathrm{Gyr}\) this is deep adiabatic, in
  which a sharp sub-orbit pericentre flip \emph{freezes out} under the
  low-pass reading (the same correction already forced on Channel A).
  Only the un-anchored \(0.45\ \mathrm{Gyr}\) value made the flip a
  resolvable dated transient.
\end{enumerate}

Consequently the post-pericentre and backsplash zones are
timescale-hostage and \textbf{not} pre-registrable; the
ancient/virialized zone is \(\sim\)zero. Only the existence claim (Tier
i, field-sector) and the first-infall-hotter leading sign (Tier ii)
survive. The corrected statement is \emph{cleaner} than the published
one: a field-sector existence claim plus a structural, timescale-free
leading sign of the opposite polarity. This paper reports the correction
with the same weight as the original claim.

\subsubsection{5.4 Magnitude and its
hostages}\label{magnitude-and-its-hostages}

At the framework-committed E10 memory the MG-impossible history spread
in the fixed-field \(\ln(\sigma_\mathrm{int}/\sigma_\mathrm{bary})\) contrast
is genuinely uncomputed \emph{within the committed kernel}: the low-pass
reading gives \(\sim1.3\%\) (canonical) / \(\sim1.5\%\) (alt), while the
framework's committed pure-phase (\(|K|=1\)) reading of the same memory
gives up to \(\sim8\)--\(13\%\) (a bounded group-delay transient
survives). We quote the magnitude as this \(\sim1.3\)--\(13\%\) band and
do not present the \(1.3\%\) low-pass corner as ``the answer'' --- it is
one corner. Channel A is \(\sim0.2\)--\(1\%\). Both are kernel-hostage:
\(\theta(y)\), the kernel \emph{shape}, and the Jensen-gap curvature are
not derived by the de Sitter--Unruh foundation, which fixes only the
cone endpoints. The instantaneous \(6\)--\(13\%\) EFE boost is a
current-configuration quantity that MG partly shares, not the
discriminant. Both footings are materially identical at fixed
dimensionless depth (\(<2\%\) on every sign fraction;
\(\tau_\mathrm{mem}H_\Lambda=2Z\) is footing-free). The leading sign is
conditional on the \(s=-1\) postulate: \(s=+1\) reverses it.

\begin{center}\rule{0.5\linewidth}{0.5pt}\end{center}

\subsection{6. Power and data
availability}\label{power-and-data-availability}

Both channels are underpowered today. We state the walls honestly and
identify what would clear them.

\subsubsection{6.1 Channel A --- star-orbit spread (ELT-gated, plus a
missing per-star
tag)}\label{channel-a-star-orbit-spread-elt-gated-plus-a-missing-per-star-tag}

The efficient (score/MLE) test has \(z=f\,w\,D\,\sqrt{2N}\), with \(f\)
the fractional \(\sigma\)-spread, \(w=\sigma^2/(\sigma^2+e^2)\) the
measurement-error down-weight (variance information falls
\emph{quadratically} in error), \(D\) the correlation of the orbit-tag
proxy with true eccentricity, and \(N\) the star count. Two independent
walls, either fatal:

\begin{itemize}
\tightlist
\item
  \textbf{Count.} At the honest \(f\sim0.2\)--\(1\%\), the Fisher floor
  needs \(N\sim7\times10^4\) (1\% ceiling) to \(\sim6\times10^5\) (0.2\%
  fiducial) clean per-star velocities in a \emph{single} deep-MOND
  system even with a perfect tag. The biggest dSph (Fornax) has
  \(\sim2600\); the whole stacked classical+diffuse reservoir is
  \(\sim7100\) (perfect-tag \(z<1.1\) at the ceiling, \(<0.4\)
  fiducial). The amplitude and the count pull \emph{opposite} --- the
  deepest-\(y\) systems (Crater II, Antlia II) that maximize \(f\) have
  the fewest stars (150--200).
\item
  \textbf{Tag.} There is no per-star eccentricity tag where the counts
  live. Gaia per-star internal proper motion has S/N
  \(\sim0.03\)--\(0.05\) (bulk systemic only, \(D_\mathrm{Gaia}\approx0\));
  LOS-only DF inference reaches \(D\lesssim0.2\) and \emph{is} the
  \(\beta\)-anisotropy channel MG reproduces; only HST/JWST multi-epoch
  3D reaches \(D\approx0.3\)--\(0.4\) for a few hundred bright stars in
  2--3 systems, giving a best real single-system \(z\approx0.05\).
\end{itemize}

No existing dataset (Walker+2009, Gaia DR3, MaNGA, ATLAS3D, Coma) comes
within 1--6 orders. Powering Channel A needs
\(\sim10^{4.5}\)--\(10^{5.5}\) clean per-star velocities in one diffuse
deep-MOND dSph from a 30 m-class campaign (ELT/MICADO, MSE) \emph{plus}
a per-star 3D orbit tag from multi-epoch space astrometry beyond Gaia's
per-star precision --- only the \(\sim1\%\) point-mass-ceiling corner is
within \(\sim1\)--\(2\) orders of plausibility. Both footings shift
\(f\) by \(<20\%\) and \(N_{3\sigma}\) by \(<44\%\): the discriminator
is magnitude- and tag-hostage, not footing-hostage.

\subsubsection{6.2 Channel B --- cluster-member EFE (scaffolding ready,
carrier count the
wall)}\label{channel-b-cluster-member-efe-scaffolding-ready-carrier-count-the-wall}

The membership and phase-tagging scaffolding is \emph{abundant and
ready}, and is not the bottleneck: GalWCat19 (Abdullah et al.~2020, ApJS
\textbf{246}:2; 1800 clusters, 34,471 members, caustic \(M/R\) at
\(\Delta=500/200/100\)), HeCS / HeCS-omnibus (Rines et al.~2013, ApJ
\textbf{767}:15; 2016, ApJ \textbf{819}:63; dedicated caustic profiles),
Yang SDSS groups, and the SAMI cluster survey (Owers et al.~2017, MNRAS
\textbf{468}:1824; 8 clusters, 2899 members). DESI DR1 BGS deepens
membership \(\sim9.5\times\) (reaching into the dwarf regime at Coma
redshift), but delivers redshifts, not IFU dispersions. The resolved
stellar \(\sigma\) is also served --- the MaNGA data-analysis-pipeline
maps recover \(\sigma_\mathrm{int}\) inside \(R_e\) (the STELLAR\_SIGMA
extensions, quadrature-subtracting the LSF correction) with no pPXF
re-derivation. The kinematics and the tagging are both fine.

The wall is the \textbf{diffuse IFU \(\sigma\)-carrier count}. The
carriers are dwarf members with reliable
\(\sigma\sim20\)--\(70\ \mathrm{km\,s^{-1}}\); MaNGA is
stellar-mass-limited (\(M_\ast\) floor \(\sim5\times10^8\)),
field-dominated, and its LSF
(\(1\sigma\approx70\)--\(76\ \mathrm{km\,s^{-1}}\)) makes reliable
stellar \(\sigma\) bottom out near \(35\)--\(45\ \mathrm{km\,s^{-1}}\),
so the \(\sigma\,20\)--\(40\) dwarfs are upper-limit-only. Worse, the
accessible \(\sim45\)--\(70\ \mathrm{km\,s^{-1}}\) carriers sit at
shallower internal depth and carry a \emph{smaller} true
\(D(\text{zone})\) than the headline band --- a tightening toward NO-GO.
Chaining the cuts (rich-cluster member \(\cap\) caustic-able \(\cap\)
Rhee-taggable \(\cap\) LSF-reliable) leaves only:

{\def\LTcaptype{none} % do not increment counter
\begin{longtable}[]{@{}
  >{\raggedright\arraybackslash}p{(\linewidth - 6\tabcolsep) * \real{0.2500}}
  >{\raggedright\arraybackslash}p{(\linewidth - 6\tabcolsep) * \real{0.2500}}
  >{\raggedright\arraybackslash}p{(\linewidth - 6\tabcolsep) * \real{0.2500}}
  >{\raggedright\arraybackslash}p{(\linewidth - 6\tabcolsep) * \real{0.2500}}@{}}
\toprule\noalign{}
\begin{minipage}[b]{\linewidth}\raggedright
configuration
\end{minipage} & \begin{minipage}[b]{\linewidth}\raggedright
diffuse tagged carriers \(N\)
\end{minipage} & \begin{minipage}[b]{\linewidth}\raggedright
exploratory S/N \(z\) (opt / mid / pes)
\end{minipage} & \begin{minipage}[b]{\linewidth}\raggedright
call
\end{minipage} \\
\midrule\noalign{}
\endhead
\bottomrule\noalign{}
\endlastfoot
public MMU-MaNGA alone & \(\sim40\)--\(77\) & \(1.2\) / \(0.5\) /
\(0.2\) & NO-GO --- dies on statistics \\
MaNGA + public SAMI stack & \(\sim135\)--\(237\) & \(2.0\) / \(0.9\) /
\(0.3\) & underpowered, firewalled hint \\
clean-exploratory floor & \(300\)--\(500\) & \(2.3\) / \(0.9\) / \(0.3\)
& required floor; still marginal \\
\end{longtable}
}

The public-MMU verdict is a \textbf{NO-GO on statistics}: the honest
central expectation is a non-detection at every accessible \(N\). Two
cautions compound it. First, the optimistic \(\sim2\)--\(2.5\sigma\)
corner assumes a \emph{signal of \(\sim9.5\%\)} (the top of the kernel
band) \emph{together with} best-case purity, scatter, and systematics;
at the low-pass \(\sim1.3\)--\(1.5\%\) signal --- comparable to or below
the \(\sim1\)--\(2\%\) projection and \(\sim2\)--\(8\%\) tidal floors
(Sec. 4.4) --- even the optimistic corner collapses and the honest
expectation is the mid column (\(z<1\)). Second, the test is
\emph{statistics}-limited at all accessible \(N\) (the statistical error
exceeds the \(\sim2\%\) systematic floor until \(N\gtrsim400\)) --- the
\emph{opposite} failure mode from an SDSS single-fiber stack, which is
\emph{systematics}-limited (its \(\sim1\%\) signal sits under a
\(1\)--\(5\%\) single-fiber \(\sigma\)-systematic floor and a
\(2\)--\(8\%\) same-signed confound, and its
\(\sigma\gtrsim90\ \mathrm{km\,s^{-1}}\) reliability floor excludes the
carriers outright). Buying \(N\) from SDSS does not buy the missing
\(\sigma\) control.

A clean detection needs \emph{either} a dedicated wide nearby-cluster
dwarf-IFU survey (with an \(M_\ast\) floor to \(\log M\sim8\), resolved
stellar \(\sigma\) reliable well below \(45\ \mathrm{km\,s^{-1}}\),
sub-percent systematics, and \(\sim10^3\)--\(10^4\) diffuse members)
\emph{or} ELT/HARMONI-class IFU (\(\sim2032\)) to reach the
\(\sim20\ \mathrm{km\,s^{-1}}\) diffuse cluster-dwarf regime. The
scaffolding slots in the moment such a carrier sample exists.

\begin{center}\rule{0.5\linewidth}{0.5pt}\end{center}

\subsection{7. Scope and limits}\label{scope-and-limits}

We restate the boundaries plainly, because they are the difference
between an honest prediction and an overclaim.

\begin{itemize}
\tightlist
\item
  \textbf{MI-class-generic, not framework-specific.} The signature
  distinguishes \emph{any} history-dependent inertia from modified
  gravity (field-sector MG \(=0\)). It does \textbf{not} distinguish
  this framework from Milgrom's linear no-EFE modified inertia
  (arXiv:2503.07106), which also produces a fixed-field spread. It is an
  MI-vs-MG test, not a this-framework-vs-Milgrom test. The
  framework-specific content is the \emph{numbers} (the committed
  \(\tau_\mathrm{mem}=2Z/H_\Lambda\), the magnitude band), not the
  discriminating power.
\item
  \textbf{It does not test \(a_0\)'s value or the sign \(s=-1\).} Both
  are postulates. At fixed dimensionless depth \(a_0\) cancels in the
  spread, so the test is \(a_0\)-value-blind by construction; the
  leading sign rides on \(s=-1\) (\(s=+1\) reverses it), so Tier (ii) is
  postulate-conditional while Tier (i) is not.
\item
  \textbf{Magnitude is kernel-hostage.} \(\theta(y)\) and the kernel
  shape (Channel B) and the Jensen-gap curvature (Channel A) are not
  derived by the de Sitter--Unruh foundation, which fixes only the cone
  endpoints. At the committed memory the Channel-B fixed-field contrast
  spans \(\sim1.3\%\) (low-pass) to \(\sim8\)--\(13\%\) (pure-delay) ---
  genuinely uncomputed \emph{within} the committed kernel. The
  \(6\)--\(13\%\) instantaneous EFE boost is a current-configuration
  quantity MG partly shares, not the discriminant. Existence (field
  sector), leading sign, and MG \(=0\) are the load-bearing claims; the
  amplitude is a band.
\item
  \textbf{MG \(=0\) is a theorem only in the field sector at fixed
  \emph{true} field.} It is \emph{not} a claim that the total observed
  spread is zero: non-equilibrium dynamics, tidal shocking, and
  projection produce same-signed fixed-field spreads in
  Newton+dark-matter and MG alike. In projection the field-sector zero
  becomes a mitigation-dependent baseline (Sec. 4.4): the raw alias is
  \(\sim1\)--\(2\%\) at realistic tag purity and can reach \(\sim7\%\)
  under filamentary/triaxial infall. The mitigation chain is
  load-bearing.
\item
  \textbf{Same-signed confounds are real and partly unmodeled.} Tidal
  heating, ram-pressure, quenching, and non-equilibrium substructure
  carry the MI sign; a radius-correlated non-equilibrium population mix
  can manufacture a large MI-signed false spread
  (\(d\ln M\approx-0.71\), \(\sim35\times\) the signal, in the
  star-orbit statistic). They are separated (in design) by the
  outward-rising radial profile (F3) and the leading sign (F2), not by
  amplitude --- and the separation is a design, not yet a demonstration
  on real data.
\item
  \textbf{Both footings throughout.} Every dimensional number is carried
  at \(a_0=9.36\times10^{-11}\) (canonical) and
  \(1.13\times10^{-10}\ \mathrm{m\,s^{-2}}\) (alt); the spread is
  footing-invariant at fixed depth and
  \(\tau_\mathrm{mem}H_\Lambda=2Z=11.58\) is footing-free.
\end{itemize}

No claim of proof is made for the framework. The field-sector MG \(=0\)
statement (at fixed true field) is the sole theorem-grade claim.

\begin{center}\rule{0.5\linewidth}{0.5pt}\end{center}

\subsection{8. Conclusion}\label{conclusion}

Post-lensing-no-go --- with lensing now shared between the MI and MG
realizations --- the relational velocity-dispersion spread is the
correct \emph{kind} of test the modified-inertia program has been
missing: an observable with an exact field-sector baseline that resists
equilibrium anisotropy. Its logical core is airtight where we claim it.
For any theory that sources a field and moves matter on WEP geodesics of
it, a subsystem's internal dispersion at fixed external field carries
exactly zero \emph{field-sector} spread across infall history; a nonzero
field-sector fixed-field history spread requires that inertia depend on
the body's own worldline, which is the definition of modified inertia.
The sole evasion --- a velocity-dependent coupling to the tracer's own
motion --- \emph{is} modified inertia in an MG costume and cannot serve
as a rival explanation. The discriminant survives a strong-anisotropy
control test in the star-orbit channel (immune by sign, the zero-point
being signal-sized), so the MI sign is not re-labeled velocity
anisotropy, and the designed observable \(D(\text{zone})\) differences
the shared radial EFE gradient out as a common mode. What the theorem
does \emph{not} deliver by itself is the \emph{observable} existence
claim: a measured fixed-field spread is MG-impossible only after the
same-signed non-field confounds --- non-equilibrium dynamics, tidal
shocking, projection --- are separated by the F3 outward-rising radial
profile. That separation is the design, not the theorem.

The framework-specific prediction, corrected from the earlier D3, is
two-tiered: the \emph{existence} of the field-sector fixed-field history
spread is theorem-grade and MG-impossible; its \emph{leading sign}, in
the first-infall pre-pericentre zone, is that first-infall members run
hotter than matched long-resident members --- a structural,
timescale-free result (for positive-averaging or pure-delay kernels)
that inverts the backwards, dated pericentre sign-flip now retracted.
The magnitude is kernel-hostage: at the framework's committed memory it
is genuinely uncomputed between a \(\sim1.3\%\) low-pass corner and a
\(\sim8\)--\(13\%\) pure-delay corner in the cluster channel, and
sub-percent in the star-orbit channel.

The honest status is a prediction, not a confrontation, and after the
magnitude correction it is a notch more cautious than a first reading
suggests. Both channels are underpowered today: the star-orbit channel
needs ELT-class per-star velocities plus a per-star 3D orbit tag that
does not exist; the cluster channel has its membership and phase-tagging
scaffolding ready (GalWCat19, HeCS, Rhee zones, DESI-deepened
membership) but is walled by the diffuse IFU \(\sigma\)-carrier count
--- a public-MMU NO-GO at \(\sim40\)--\(77\) carriers, an underpowered
\(\sim2\)--\(2.5\sigma\) firewalled hint with a MaNGA+SAMI stack (and
only if the signal sits at the top of the kernel band; at the low-pass
corner even that collapses to a non-detection), and a clean detection
gated on a dedicated wide-cluster dwarf-IFU survey or ELT/HARMONI
(\(\sim2032\)). This front is the program's most distinctive
\emph{in-principle} discriminant, worth pre-registering (field-sector
existence plus first-infall-hotter, with the sign statistic pinned to
the single robust zone) and reanalysing at MaNGA/SAMI for a hint --- but
it neither manufactures a win nor a deficit, and it does not test
\(a_0\)'s value or the sign postulate.

\begin{center}\rule{0.5\linewidth}{0.5pt}\end{center}

\subsection{Appendix A --- Committed, verifiable
scripts}\label{appendix-a-committed-verifiable-scripts}

Every load-bearing claim is backed by a committed, runnable script
(numpy/scipy/sympy), each exiting 0 and carrying both coefficient
footings. The four lanes:

\textbf{Star-orbit channel} (\texttt{prep\_2026/sigma\_spread/}): -
\texttt{mi\_spread.py} --- the deep-adiabatic Jensen-gap magnitude
(\(\tau_\mathrm{mem}/T_\mathrm{orb}\) table; RMS \(\sigma\)-spread
\(0.2\)--\(1\%\); sign negative), cross-checked against the
19/19-verified real-kernel integrator
(\texttt{prep\_2026/mi\_integrator/}). - \texttt{mg\_zero.py} --- the
sourced-field--WEP MG \(=0\) theorem (\(d\sigma/dy=0\) symbolic;
boundary stress test C1--C6; the disformal/Finsler-SME evasion
identified as MI). - \texttt{observable.py} --- the \(\beta\)-immune
orbit-family enclosed-mass consistency statistic, the anisotropy
sign-immunity control (zero-point \(+0.033\), signal-sized, immune by
sign), and the non-equilibrium population-mix false-detection probe
(\(d\ln M\approx-0.71\)). - \texttt{power.py} --- the Fisher floor,
MC-validated (\(\sqrt N\) scaling \(1.99\); null calibrated), and the
two-wall no-go.

\textbf{Cluster-member EFE channel}
(\texttt{prep\_2026/cluster\_efe\_channel/}): - \texttt{predict.py} ---
the \(\theta(y)\) magnitude band, radial structure, and mass/depth
dependence. - \texttt{mg\_efe\_zero.py} --- the fixed-true-field
field-sector MG \(=0\) theorem and the quantified projection/interloper
mimics. - \texttt{observable.py} --- the \(D(\text{zone})\) statistic
and the four-part confound fingerprint. - \texttt{power.py} --- the
survey-regime power analysis (SDSS systematics-limited; MaNGA/SAMI
statistics-limited exploratory).

\textbf{Sign reconciliation} (\texttt{prep\_2026/cluster\_efe\_sign/}):
- \texttt{setup\_diagnose.py}, \texttt{net\_sign.py},
\texttt{robustness.py} --- the diagnosis of the two label bugs on one
correct baseline, the per-zone net sign against both anchors (low-pass
and pure-delay kernel readings), the timescale pin, and the \(0/125\)
orbit-scan robustness of the first-infall-hotter sign.

\textbf{Data-availability scout}
(\texttt{prep\_2026/cluster\_efe\_data\_scout/}): -
\texttt{inventory\_manga.py}, \texttt{inventory\_membership.py},
\texttt{overlap\_power.py} --- the MMU carrier census (resolved
\(\sigma\) served; membership abundant; carrier count \(\sim40\)--\(77\)
the binding wall).

\begin{center}\rule{0.5\linewidth}{0.5pt}\end{center}

\subsection{References}\label{references}

\begin{itemize}
\tightlist
\item
  Abdullah, M. H., Wilson, G., Klypin, A., et al.~2020, ApJS
  \textbf{246}, 2 (the GalWCat19 cluster catalog).
\item
  Milgrom, M. 1983, ApJ \textbf{270}, 365.
\item
  Milgrom, M. 1999, Phys. Lett. A \textbf{253}, 273 (astro-ph/9805346)
  --- the \(\nu\)-kernel wellhead, \(\nu(y)=\sqrt{1+1/y}\), Eq. 9.
\item
  Milgrom, M. 2014, MNRAS \textbf{437}, 2531 --- the MOND paradigm; the
  external field effect and MG-virial universality.
\item
  Milgrom, M. 2022, Phys. Rev.~D \textbf{106}, 064060 (cited for the
  two-frequency EFE subsystem boost \(\theta(y)\)).
\item
  Milgrom, M. 2025, arXiv:2503.07106 --- a linear no-EFE
  modified-inertia model (also produces a fixed-field spread; the
  MI-class-generic caveat).
\item
  Owers, M. S., et al.~2017, MNRAS \textbf{468}, 1824 (the SAMI cluster
  redshift survey).
\item
  Rhee, J., Smith, R., Choi, H., et al.~2017, ApJ \textbf{843}, 128 ---
  projected-phase-space infall zones.
\item
  Rines, K., et al.~2013, ApJ \textbf{767}, 15; 2016, ApJ \textbf{819},
  63 --- the HeCS / HeCS-omnibus caustic cluster profiles.
\item
  Wolf, J., Martinez, G. D., Bullock, J. S., et al.~2010, MNRAS
  \textbf{406}, 1220 --- the \(\beta\)-immune half-light mass estimator.
\item
  SPARC / dSph kinematics: Walker et al.~2009 (ApJ \textbf{704}, 1274)
  and the Battaglia et al.~/ Gaia dSph proper-motion literature, as
  needed.
\item
  Zimmerman, C. 2026, MI Field Theory Results 2026, Zenodo
  {[}10.5281/zenodo.21403470{]}.
\item
  Zimmerman, C. 2026, the lensing no-go, Zenodo
  {[}10.5281/zenodo.21418816{]}.
\item
  Zimmerman, C. 2026, No Pump-Free Corner (superseded on D3), Zenodo
  {[}10.5281/zenodo.21179352{]}, with the concurrent erratum
  (2026-07-17).
\end{itemize}

\emph{\(a_0\)'s value and the sign \(s=-1\) remain postulates.
Field-sector MG \(=0\) at fixed true field is the sole theorem-grade
claim. No claim of proof is made for the framework. Both coefficient
footings (\(a_0=9.36\times10^{-11}\) /
\(1.13\times10^{-10}\ \mathrm{m\,s^{-2}}\)) are carried throughout.}

\end{document}
